\section{Data}
\label{sec:data}

\subsection{Pools}

Each pool pairs a single receptor with a molecular library scored by one
docking program, so every row is one molecule and the label is its docking
score (more negative indicates a better predicted binder). Because the receptor
is fixed within a pool, no protein representation enters the model: the task is
a regression on molecular representation alone.

\begin{table}[t]
\centering
\caption{Docking pools. Scores are not comparable in absolute terms across
oracles; recall and enrichment are.}
\label{tab:pools}
\small
\begin{tabular}{llrrl}
\toprule
Pool & Target & Molecules & Backbones & Oracle \\
\midrule
\texttt{enhits} & EnamineHTS                & \num{2104318}    & 5 & DOCK3.7 \\
\texttt{ampc}   & AmpC $\beta$-lactamase    & \num{98489350}   & 6 & DOCK3.7 \\
\texttt{d4}     & D4 dopamine receptor      & \num{116225927}  & 6 & DOCK3.7 \\
\texttt{sura}   & SurA chaperone            & \num{1531875935} & \todo{--} & Glide HTVS \\
\bottomrule
\end{tabular}
\end{table}

The \texttt{sura} pool is listed for completeness but excluded from all
experiments reported here: its embeddings exist only as partial, unstitched
per-chunk arrays at the time of writing. Its inclusion matters for a future
version of this work precisely because it uses a different docking program;
until then, every result below is a statement about DOCK3.7 specifically rather
than about docking in general.

\subsection{Representations}

Each pool is accompanied by precomputed embeddings from several pretrained
molecular models, stored as dense \texttt{float32} arrays row-aligned with the
molecule list (Table~\ref{tab:backbones}). Storage is a material constraint: a
single \num{3400}-dimensional embedding of the D4 pool occupies
\SI{1.58}{\tera\byte}.

\begin{table}[t]
\centering
\caption{Embedding backbones, dimensionality, and storage for the D4 pool
(\num{116225927} molecules, \texttt{float32}).}
\label{tab:backbones}
\small
\begin{tabular}{llrr}
\toprule
Backbone & Modality & Dim. & D4 storage (\si{\tera\byte}) \\
\midrule
GROVER~\citep{rong2020grover}        & molecular graph     & 3400 & 1.58 \\
Uni-Mol2~\citep{zhou2023unimol}      & 3D conformer        & 1536 & 0.71 \\
MHG-GED~\citep{kajino2019mhg}        & hypergraph grammar  & 1024 & 0.48 \\
MolFormer~\citep{ross2022molformer}  & SMILES              &  768 & 0.36 \\
SMI-TED~\citep{soares2025smited}     & SMILES              &  768 & 0.36 \\
Uni-Mol v1~\citep{zhou2023unimol}    & 3D conformer        &  512 & 0.24 \\
\bottomrule
\end{tabular}
\end{table}

Two inconsistencies in the stored artifacts are worth recording, as they are
invisible to a consumer who assumes uniformity. On \texttt{enhits}, the
Uni-Mol2 embedding is \num{768}-dimensional rather than \num{1536}, having been
produced by a different extraction; results for that backbone are consequently
not comparable across pools. The same pool also lacks Uni-Mol v1 entirely and
stores GROVER under a different filename convention.

\subsection{Row alignment}
\label{sec:alignment}

Surrogate benchmarking of this kind depends on an assumption that is easy to
state and easy to violate: that row $i$ of the embedding array corresponds to
row $i$ of the score vector. A violation does not raise an error. It produces a
model that trains to a plausible-looking loss curve and reports meaningless
numbers.

For two of the three pools used here the assumption holds directly. For
\texttt{enhits} it does not: the embeddings are stored in a shuffled order
relative to the only available score file, which is itself sorted by score, so
a naive join yields labels that are both permuted and systematically ordered.
We therefore treat alignment as a verification step rather than an assumption.
Our pipeline reconstructs the permutation, applies it, and then checks the
molecular identifier of every row against a reference list, refusing to proceed
on any mismatch. All \num{2104318} rows of \texttt{enhits} pass this check; we
additionally confirmed agreement on \num{50000} randomly drawn rows through an
independent identifier-to-score mapping.

We report this in detail because it is a reproducibility hazard rather than an
incidental engineering note: any benchmark built on these public artifacts
without such a check may silently be reporting noise.

\subsection{Failed docking runs}

Docking runs that fail to produce a pose, or that produce severe steric
clashes, yield non-negative scores that are not on the same scale as successful
results. These are removed before splitting, so that training, validation and
test sets share one universe and the filter is recorded alongside the split
rather than applied ad hoc at training time. The affected fraction varies
sharply by pool: approximately \SI{5}{\percent} of the D4 pool (maximum
observed score $+328$), a negligible fraction of AmpC, and \num{16} rows of
\texttt{enhits}.
